\section{CPU 如何执行指令}

\subsection{本章要解决的问题}

上一章已经把一份 C/C++ 源码拆成了预处理、编译、汇编、链接、加载和进程执行。那条路径告诉我们：源代码最后会变成可执行文件中的机器码字节，操作系统把这些字节映射到进程的虚拟地址空间里，然后 CPU 从某个入口地址开始执行。

本章继续往下一层走。我们要回答一个看似简单、实际非常关键的问题：

\begin{center}
\textbf{CPU 到底怎样把一条指令变成一次真实的机器行为？}
\end{center}

对初学者来说，这个问题通常会被几种不准确的说法遮住：

\begin{itemize}
  \item ``CPU 执行 C++ 代码''。
  \item ``汇编就是 CPU 直接看到的东西''。
  \item ``一条汇编就是一个硬件动作''。
  \item ``CPU 一条一条顺序执行，所以源代码顺序就是真实执行顺序''。
  \item ``内存就是一个大数组，访问一个变量就是直接去内存取它''。
\end{itemize}

这些说法都有一点影子，但都不够精确。为了做高性能计算、AI 算子优化和 CPU kernel 开发，必须建立更细的模型。我们不要求你在本章就理解现代乱序执行核心的所有细节，但你必须形成一条稳定的链路：

\begin{center}
C/C++ 表达式 $\rightarrow$ 汇编文本 $\rightarrow$ 机器码字节 $\rightarrow$ 指令语义 $\rightarrow$ 架构状态变化 $\rightarrow$ 微架构执行过程 $\rightarrow$ 性能现象
\end{center}

后面所有优化都会反复使用这条链路。你看一个循环为什么慢，最终还是要问：

\begin{itemize}
  \item 它生成了哪些指令？
  \item 这些指令读写哪些寄存器和内存？
  \item 哪些指令依赖前面的结果？
  \item 哪些访问会打到 cache，哪些会去更远的内存层级？
  \item 哪些分支会打断前端供给？
  \item CPU 内部能否把多条指令重叠执行？
\end{itemize}

\begin{keyidea}
本章不是为了把 CPU 讲成神秘黑盒，而是为了把 CPU 拆成几个可观察、可推理、可实验验证的层次。性能工程的第一原则是：先知道机器真正执行了什么，再讨论优化。
\end{keyidea}

\subsection{从一个最小函数开始}

考虑这个 C++ 函数：

\begin{lstlisting}[language=C++]
extern "C" int add_i32(int x, int y)
{
    return x + y;
}
\end{lstlisting}

在 x86-64 System V ABI 下，前两个整数参数通常通过 \register{edi} 和 \register{esi} 传入，32 位整数返回值通常放在 \register{eax}。一个容易理解的汇编版本可以写成：

\begin{lstlisting}
mov eax, edi
add eax, esi
ret
\end{lstlisting}

这三行已经包含了本章的核心问题：

\begin{itemize}
  \item \instruction{mov} 为什么能把一个参数复制到返回寄存器？
  \item \instruction{add} 读了什么，写了什么？
  \item \instruction{ret} 为什么能回到调用者？
  \item CPU 怎样知道下一条该执行哪条指令？
  \item 这些汇编文本进入可执行文件后到底是什么？
\end{itemize}

先强调一个事实：CPU 不知道变量名 \code{x} 和 \code{y}。变量名、函数名、类型名是源代码和调试信息层面的东西。CPU 执行时看到的是指令字节，操作的是寄存器、内存地址和状态位。

如果把这个函数放进目标文件，用 \code{objdump -d -Mintel} 观察，机器码字节可能类似：

\begin{lstlisting}
89 f8        mov eax, edi
01 f0        add eax, esi
c3           ret
\end{lstlisting}

也可能在优化后看到：

\begin{lstlisting}
8d 04 37     lea eax, [rdi+rsi]
c3           ret
\end{lstlisting}

编译器选择哪一种，取决于优化级别、目标 CPU、上下文和它是否需要保留 flags。这里先不急着判断哪种更好。本章先建立一个原则：

\begin{mentalmodel}
汇编是人类可读的指令表示，机器码是 CPU 取指单元读取的字节序列。C++ 源码不是 CPU 的输入，汇编文本通常也不是 CPU 的最终输入；真正被执行的是内存中的机器码字节。
\end{mentalmodel}

\subsection{CPU 是什么：从状态机开始理解}

很多教材会直接说 CPU 由控制器、运算器和寄存器组成。这个说法没错，但对性能学习还不够。我们需要更底层一点的心智模型：CPU 是一个巨大而高速的状态机。

所谓状态机，就是它有一组状态，每个时钟周期根据当前状态和输入计算出下一个状态。对程序员来说，最重要的状态包括：

\begin{itemize}
  \item 当前要执行的位置，也就是 instruction pointer，在 x86-64 中叫 \register{rip}。
  \item 通用寄存器，例如 \register{rax}、\register{rbx}、\register{rcx}、\register{rdx}、\register{rdi}、\register{rsi}。
  \item 栈指针 \register{rsp}。
  \item 条件标志位 RFLAGS，例如 zero flag、carry flag、sign flag、overflow flag。
  \item 内存中可见的数据。
\end{itemize}

从硬件实现角度看，CPU 内部有大量组合逻辑和时序逻辑。

\term{组合逻辑}{combinational logic} 的输出只由当前输入决定。加法器、比较器、多路选择器、地址生成单元都可以从这个角度理解。给加法器两个输入位模式，它经过门电路传播后给出结果和进位。

\term{时序逻辑}{sequential logic} 会保存状态。寄存器、流水线寄存器、重排序缓冲区中的条目、cache tag 阵列中的状态都属于这一类。时钟边沿到来时，一部分计算结果被锁存成新的状态。

你不需要在本书前几章就掌握电路设计，但必须理解一点：机器执行不是抽象概念，它最终落实为状态的读、组合逻辑的计算和状态的写。

\begin{center}
\begin{minipage}{0.82\linewidth}
\begin{verbatim}
old state + instruction bytes
        |
        v
decode instruction meaning
        |
        v
read operands -> compute -> write results
        |
        v
new state
\end{verbatim}
\end{minipage}
\end{center}

这就是最朴素的 CPU 执行图景。现代 CPU 会把这个过程拆开、重叠、预测、乱序执行、再按程序顺序提交结果，但它必须最终制造出和 ISA 规定一致的可见状态。

\subsection{ISA：软件和硬件之间的契约}

\term{ISA}{instruction set architecture，指令集体系结构} 是软件和处理器之间的契约。对 x86-64 来说，ISA 规定了：

\begin{itemize}
  \item 有哪些寄存器。
  \item 每条指令的编码格式。
  \item 每条指令读写什么操作数。
  \item 指令如何改变寄存器、内存、flags 和 \register{rip}。
  \item 异常、中断、权限、页表等系统级行为如何表现。
\end{itemize}

ISA 不规定 CPU 内部必须如何实现。例如同样是执行 \instruction{add eax, esi}：

\begin{itemize}
  \item 一个简单教学 CPU 可以顺序取指、解码、读寄存器、加法、写回。
  \item 一个现代桌面 CPU 可能把它解码为 micro-op，进入调度队列，等待操作数 ready，在某个整数执行端口上执行，最后在 retire 阶段提交。
  \item 一个仿真器可以用软件解释这条指令。
\end{itemize}

只要对程序可见的结果一致，它们都符合 ISA。

\term{架构状态}{architectural state} 是 ISA 规定的软件可见状态，例如通用寄存器、\register{rip}、RFLAGS 和内存。 \term{微架构}{microarchitecture} 是某个具体 CPU 实现 ISA 的内部方式，例如 pipeline、decode queue、uop cache、register renaming、reservation station、load/store queue、L1 cache、TLB 和 branch predictor。

\begin{keyidea}
ISA 决定正确性语义，微架构决定性能表现。同一份 x86-64 程序在不同 CPU 上结果应该一致，但速度可能差很多，因为内部实现不同。
\end{keyidea}

\subsection{x86-64 程序员模型}

学习汇编时，先不要把 CPU 想成无限复杂的乱序机器。先从程序员模型开始。所谓程序员模型，就是写汇编、读反汇编、理解 ABI 时必须看到的那部分 CPU 状态。

\subsubsection{通用寄存器}

x86-64 有 16 个常用通用整数寄存器：

\begin{center}
\begin{tabular}{lll}
\toprule
64 位 & 32 位低半部分 & 常见用途 \\
\midrule
\register{rax} & \register{eax} & 返回值、通用计算 \\
\register{rbx} & \register{ebx} & callee-saved 通用寄存器 \\
\register{rcx} & \register{ecx} & 第 4 个整数参数、计数、通用计算 \\
\register{rdx} & \register{edx} & 第 3 个整数参数、乘除辅助、通用计算 \\
\register{rsi} & \register{esi} & 第 2 个整数/指针参数 \\
\register{rdi} & \register{edi} & 第 1 个整数/指针参数 \\
\register{rbp} & \register{ebp} & 栈帧指针或通用 callee-saved 寄存器 \\
\register{rsp} & \register{esp} & 栈指针 \\
\register{r8}  & \register{r8d} & 第 5 个整数/指针参数 \\
\register{r9}  & \register{r9d} & 第 6 个整数/指针参数 \\
\register{r10} & \register{r10d} & caller-saved 临时寄存器 \\
\register{r11} & \register{r11d} & caller-saved 临时寄存器 \\
\register{r12} & \register{r12d} & callee-saved 通用寄存器 \\
\register{r13} & \register{r13d} & callee-saved 通用寄存器 \\
\register{r14} & \register{r14d} & callee-saved 通用寄存器 \\
\register{r15} & \register{r15d} & callee-saved 通用寄存器 \\
\bottomrule
\end{tabular}
\end{center}

寄存器不是变量。寄存器是 CPU 内部的高速状态存储单元。编译器会把源代码里的变量、临时值、地址、中间计算结果分配到寄存器或栈位置。一个变量可能根本没有固定的寄存器；一个寄存器在不同指令之间也可能承载不同含义。

\subsubsection{32 位写入会清零高 32 位}

x86-64 有一个非常重要的规则：写入 32 位寄存器通常会把对应 64 位寄存器的高 32 位清零。例如：

\begin{lstlisting}
mov eax, edi
\end{lstlisting}

这条指令写 \register{eax}，同时会让 \register{rax} 的高 32 位变成 0。这个规则经常帮助 CPU 避免旧高位造成的依赖，也让 32 位整数计算自然产生零扩展结果。

\begin{warningbox}
不要把 \register{eax} 理解成完全独立于 \register{rax} 的寄存器。它是 \register{rax} 的低 32 位视图。类似地，\register{ax}、\register{al} 也是同一个物理架构寄存器的更小视图。部分寄存器写入在某些老架构上还会带来额外依赖问题，后面讲性能时会再次出现。
\end{warningbox}

\subsubsection{RIP：CPU 下一步从哪里取指}

\register{rip} 是 instruction pointer。它保存下一条要执行的指令地址。普通顺序执行时，CPU 取出当前指令后，会把 \register{rip} 更新到下一条指令地址。

x86 指令长度不是固定的。刚才的例子中：

\begin{lstlisting}
89 f8        mov eax, edi
01 f0        add eax, esi
c3           ret
\end{lstlisting}

\instruction{mov eax, edi} 长 2 字节，\instruction{add eax, esi} 长 2 字节，\instruction{ret} 长 1 字节。所以如果第一条指令地址是 \code{0x1000}，顺序取指时大致会这样：

\begin{center}
\begin{tabular}{lll}
\toprule
执行前 \register{rip} & 指令字节 & 顺序下一地址 \\
\midrule
\code{0x1000} & \code{89 f8} & \code{0x1002} \\
\code{0x1002} & \code{01 f0} & \code{0x1004} \\
\code{0x1004} & \code{c3}    & 由栈中的返回地址决定 \\
\bottomrule
\end{tabular}
\end{center}

跳转、调用、返回和异常都会改变 \register{rip}。控制流的本质就是改变 CPU 接下来从哪里取指。

\subsubsection{RFLAGS：条件判断的隐含状态}

很多整数指令会设置 RFLAGS。例如 \instruction{add} 不只写目标寄存器，还会更新 zero flag、sign flag、carry flag、overflow flag 等状态。条件跳转指令读取这些 flags 来决定是否跳转。

例如：

\begin{lstlisting}
cmp edi, esi
jg  greater
\end{lstlisting}

\instruction{cmp edi, esi} 的语义类似计算 \code{edi - esi} 并更新 flags，但不保存结果。 \instruction{jg} 读取 flags，判断有符号意义下 \register{edi} 是否大于 \register{esi}。

这解释了为什么 \instruction{lea eax, [rdi+rsi]} 有时会替代 \instruction{add}：\instruction{lea} 做地址形式的整数计算，但通常不改 flags。如果后续代码还需要旧 flags，使用 \instruction{lea} 可以避免破坏它们；如果后续不需要 flags，编译器也可能因为编码或调度原因选择其中一种。

\subsection{机器码字节和汇编文本}

汇编文本是给人看的。机器码字节才是 CPU 取指单元读取的内容。汇编器负责把：

\begin{lstlisting}
mov eax, edi
\end{lstlisting}

变成：

\begin{lstlisting}
89 f8
\end{lstlisting}

这两个字节不是随意来的。x86 指令编码中包含 opcode、寄存器编号、寻址模式、立即数、位移、前缀等部分。完整 x86 编码非常复杂，本章只建立直觉。

以 \instruction{mov eax, edi} 为例：

\begin{itemize}
  \item \code{89} 是一类 \instruction{mov} 指令的 opcode。
  \item \code{f8} 是 ModRM 字节，描述源寄存器和目标寄存器。
\end{itemize}

以 \instruction{ret} 为例：

\begin{itemize}
  \item \code{c3} 就是一字节 near return 指令。
\end{itemize}

x86-64 是变长指令集。一条指令可以很短，也可以很长。变长编码带来代码密度优势，也让前端解码变复杂。现代 x86 CPU 为了高速执行，必须在前端投入大量硬件：取指、边界识别、解码、缓存已经解码过的 uops 等。后面讲前端瓶颈时，这一点会非常重要。

\begin{mentalmodel}
一条汇编指令不是一个字符串，而是一段字节的解释方式。反汇编器把字节翻译回人类可读文本；汇编器把文本翻译成字节。性能分析时，二者都要看，但最终以实际二进制中的字节和指令为准。
\end{mentalmodel}

\subsection{最朴素的 Fetch-Decode-Execute 模型}

先构造一个极简 CPU 模型。它每次只执行一条指令，每条指令完整执行完才开始下一条：

\begin{center}
\begin{minipage}{0.78\linewidth}
\begin{verbatim}
while running:
    bytes = fetch memory[RIP]
    inst = decode(bytes)
    next_RIP = RIP + length(inst)
    operands = read_registers_or_memory(inst)
    result = execute(inst, operands)
    write_registers_or_memory(result)
    RIP = choose_next_RIP(inst, next_RIP)
\end{verbatim}
\end{minipage}
\end{center}

这个模型非常简化，但它的每个词都值得解释。

\subsubsection{Fetch：取指}

取指就是根据 \register{rip} 到内存层级中取机器码字节。现代 CPU 通常不是每次直接去主内存取字节，而是从 L1 instruction cache、uop cache 或其他前端结构中拿。第一次学习时可以先想象：

\begin{center}
\register{rip} $\rightarrow$ instruction memory $\rightarrow$ instruction bytes
\end{center}

如果指令字节不在前端 cache 里，取指就可能变慢。代码太大、分支太乱、热循环跨越不友好的边界，都可能影响前端供给。

\subsubsection{Decode：解码}

解码器把机器码字节解释成指令语义。它要知道：

\begin{itemize}
  \item 这是什么指令。
  \item 指令有多长。
  \item 操作数在哪里。
  \item 是否有立即数或内存位移。
  \item 这条指令应该读写哪些架构状态。
\end{itemize}

现代 x86 CPU 通常会把复杂 x86 指令解码成更接近内部执行格式的 micro-ops，简称 uops。不要把 uops 当成另一套必须立刻精通的汇编语言；现在只需要知道：软件看到 x86 指令，硬件内部可能执行一个或多个更小的操作。

\subsubsection{Read operands：读取操作数}

指令需要输入。输入可能来自：

\begin{itemize}
  \item 寄存器，例如 \register{edi}、\register{esi}。
  \item 立即数，例如 \code{add eax, 5} 里的 \code{5}。
  \item 内存，例如 \code{mov eax, dword ptr [rdi]}。
  \item 隐含状态，例如 \instruction{ret} 隐含读取 \register{rsp} 指向的栈内存。
\end{itemize}

读寄存器通常很快。读内存则复杂得多，因为地址可能还要计算，虚拟地址可能要经过 TLB 翻译，数据可能在 L1/L2/L3 cache 或主内存。

\subsubsection{Execute：执行}

执行是根据指令语义完成计算或动作：

\begin{itemize}
  \item 整数加法由整数执行单元完成。
  \item 地址计算由地址生成单元完成。
  \item load 从内存层级读取数据。
  \item store 把数据写向内存系统。
  \item branch 计算下一条指令地址。
  \item floating-point 和 SIMD 指令使用向量/浮点执行单元。
\end{itemize}

\subsubsection{Write back / Commit：写回和提交}

计算结果最终要改变状态。简单 CPU 可以执行完就写架构寄存器。现代乱序 CPU 会先把结果写到内部物理寄存器或缓冲结构中，再在 retire 阶段按程序顺序提交。这样可以保证异常和中断时仍然呈现精确的架构状态。

本章先记住一句话：

\begin{keyidea}
CPU 可以在内部提前做很多事，但最终必须让软件看到像是按照程序顺序执行每条指令一样的结果。这是理解乱序执行和精确异常的基础。
\end{keyidea}

\subsection{逐条执行 \code{add_i32}}

现在手工执行这个版本：

\begin{lstlisting}
mov eax, edi
add eax, esi
ret
\end{lstlisting}

假设函数入口处：

\begin{itemize}
  \item \register{edi} = 10。
  \item \register{esi} = 32。
  \item \register{rsp} 指向栈顶。
  \item 栈顶 8 字节保存调用者的返回地址。
\end{itemize}

第一条：

\begin{lstlisting}
mov eax, edi
\end{lstlisting}

语义是把 \register{edi} 的低 32 位复制到 \register{eax}。执行后：

\begin{itemize}
  \item \register{eax} = 10。
  \item \register{rax} 高 32 位清零。
  \item 通常不改变 flags。
  \item \register{rip} 指向下一条指令。
\end{itemize}

第二条：

\begin{lstlisting}
add eax, esi
\end{lstlisting}

语义是：

\begin{center}
\register{eax} $\leftarrow$ \register{eax} + \register{esi}
\end{center}

执行后：

\begin{itemize}
  \item \register{eax} = 42。
  \item RFLAGS 根据结果更新。
  \item \register{rax} 高 32 位仍然是 0。
  \item \register{rip} 指向下一条指令。
\end{itemize}

第三条：

\begin{lstlisting}
ret
\end{lstlisting}

近似语义是：

\begin{center}
\begin{minipage}{0.72\linewidth}
\begin{verbatim}
target = memory[rsp]
rsp = rsp + 8
rip = target
\end{verbatim}
\end{minipage}
\end{center}

也就是说，\instruction{ret} 从栈顶弹出返回地址，更新 \register{rsp}，然后把 \register{rip} 改成返回地址。函数返回值已经在 \register{eax} 中，调用者按照 ABI 从那里读取返回值。

\subsection{调用和返回：栈为什么参与控制流}

函数调用不是魔法。假设调用者执行：

\begin{lstlisting}
call add_i32
\end{lstlisting}

\instruction{call} 的关键语义可以近似理解为：

\begin{center}
\begin{minipage}{0.72\linewidth}
\begin{verbatim}
rsp = rsp - 8
memory[rsp] = address_of_next_instruction
rip = address_of_add_i32
\end{verbatim}
\end{minipage}
\end{center}

所以 \instruction{call} 做了两件事：

\begin{itemize}
  \item 保存返回地址。
  \item 跳转到被调用函数入口。
\end{itemize}

\instruction{ret} 则做反向操作：

\begin{itemize}
  \item 从栈顶取回返回地址。
  \item 跳回调用者。
\end{itemize}

这解释了为什么破坏 \register{rsp} 会让程序崩溃，也解释了为什么栈溢出、返回地址覆盖和控制流劫持是安全问题。对性能学习来说，理解 call/ret 也很重要，因为函数调用涉及 ABI、寄存器保存、栈对齐、返回地址预测和 inline 决策。

\begin{warningbox}
\instruction{call} 和 \instruction{ret} 改变的是控制流，不是普通数据流。你必须同时追踪 \register{rip}、\register{rsp} 和栈内存，才能正确理解函数调用。
\end{warningbox}

\subsection{内存访问：一条 load 背后的路径}

再看一个稍微复杂的函数：

\begin{lstlisting}[language=C++]
extern "C" int add_load(const int* p, int y)
{
    return p[0] + y;
}
\end{lstlisting}

可能出现这样的汇编：

\begin{lstlisting}
mov eax, dword ptr [rdi]
add eax, esi
ret
\end{lstlisting}

这里 \register{rdi} 保存指针 \code{p}，\code{dword ptr [rdi]} 表示从地址 \register{rdi} 处读取 4 字节整数。

从 ISA 语义看，这条 load 很简单：

\begin{center}
\register{eax} $\leftarrow$ memory[\register{rdi} ... \register{rdi}+3]
\end{center}

从微架构看，它可能经过很多步骤：

\begin{enumerate}
  \item 地址生成单元计算有效地址。
  \item 如果使用虚拟内存，CPU 用 TLB 查询虚拟地址到物理地址的翻译。
  \item L1 data cache 根据地址查找 cache line。
  \item 如果 L1 命中，数据很快返回。
  \item 如果 L1 未命中，继续查 L2、L3，甚至访问主内存。
  \item 取回 4 字节并送到目标寄存器。
\end{enumerate}

这就是性能工程中经常说的：寄存器计算和内存访问不是一个数量级的问题。一个整数加法可能吞吐很高，而一次真正去主内存的 load 会慢很多。

\subsection{地址计算：数组访问并不神秘}

考虑：

\begin{lstlisting}[language=C++]
extern "C" int get_i(const int* a, long i)
{
    return a[i];
}
\end{lstlisting}

核心地址计算是：

\begin{center}
\code{address = a + i * sizeof(int)}
\end{center}

x86-64 寻址模式能直接表达常见的 base + index * scale + displacement：

\begin{lstlisting}
mov eax, dword ptr [rdi + rsi*4]
ret
\end{lstlisting}

这里：

\begin{itemize}
  \item \register{rdi} 是 base，也就是数组起始地址。
  \item \register{rsi} 是 index，也就是下标。
  \item \code{4} 是 scale，因为 \code{int} 通常 4 字节。
  \item 没有额外 displacement。
\end{itemize}

这条指令同时包含地址计算和内存读取。现代 CPU 内部通常会把地址生成、TLB 查询、cache 访问等步骤分到 load pipeline 中。后面学习 cache、TLB 和 prefetch 时，这个例子会继续使用。

\subsection{控制流：顺序、跳转和分支}

没有跳转时，\register{rip} 顺序前进。有跳转时，下一条指令地址由跳转目标决定。

例如：

\begin{lstlisting}[language=C++]
extern "C" int abs_i32(int x)
{
    if (x < 0) {
        return -x;
    }
    return x;
}
\end{lstlisting}

一个未必最优但容易理解的汇编形态是：

\begin{lstlisting}
mov eax, edi
test edi, edi
jge done
neg eax
done:
ret
\end{lstlisting}

\instruction{test edi, edi} 会计算 \register{edi} 和自身的按位与，只更新 flags，不保存结果。它常用于判断是否为 0 或正负。 \instruction{jge done} 根据 flags 决定是否跳到 \code{done}。

控制流引入了性能问题：CPU 前端必须知道接下来取哪里。如果等分支真正执行完再取下一条，流水线会经常停住。现代 CPU 因此引入 branch predictor，提前猜测分支方向和目标地址。

\subsection{为什么需要流水线}

假设一条指令要经历五步：

\begin{enumerate}
  \item 取指。
  \item 解码。
  \item 读操作数。
  \item 执行。
  \item 写回。
\end{enumerate}

如果一条指令完成所有步骤后才开始下一条，硬件资源大部分时间都在空闲。流水线的想法是把不同指令的不同阶段重叠起来：

\begin{center}
\begin{tabular}{llllll}
\toprule
周期 & 阶段 1 & 阶段 2 & 阶段 3 & 阶段 4 & 阶段 5 \\
\midrule
1 & I1 取指 & & & & \\
2 & I2 取指 & I1 解码 & & & \\
3 & I3 取指 & I2 解码 & I1 读数 & & \\
4 & I4 取指 & I3 解码 & I2 读数 & I1 执行 & \\
5 & I5 取指 & I4 解码 & I3 读数 & I2 执行 & I1 写回 \\
\bottomrule
\end{tabular}
\end{center}

流水线提高的是吞吐，不一定降低单条指令延迟。这里要提前区分两个性能概念：

\term{延迟}{latency} 是一个操作从开始到结果可用需要多久。 \term{吞吐}{throughput} 是单位时间内可以完成多少操作。

一条加法可能延迟 1 个周期，但如果 CPU 每周期能发射多条独立加法，那么吞吐可以高于每周期 1 条。后面讲 latency、throughput、ILP 时会严格量化这些概念。

\subsection{依赖：为什么有些指令不能并行}

考虑：

\begin{lstlisting}
add eax, esi
add eax, edx
add eax, ecx
\end{lstlisting}

每条指令都读写 \register{eax}。第二条必须等待第一条产生新的 \register{eax}，第三条必须等待第二条。这叫数据依赖。

再考虑：

\begin{lstlisting}
add eax, esi
add r8d, r9d
add r10d, r11d
\end{lstlisting}

这三条指令互相独立。现代 CPU 更容易把它们重叠执行。高性能代码经常要制造足够多的独立工作，让 CPU 的多个执行单元都有活干。

\begin{keyidea}
性能优化不是简单减少源代码行数，而是控制机器指令的依赖图、内存访问模式和前端供给。短代码不一定快，能让硬件持续高效工作的代码才快。
\end{keyidea}

\subsection{为什么需要乱序执行}

顺序流水线遇到长延迟操作会卡住。例如：

\begin{lstlisting}
mov eax, dword ptr [rdi]   ; 可能 cache miss
add r8d, r9d               ; 与 eax 无关
add r10d, r11d             ; 与 eax 无关
add eax, esi               ; 依赖 load
\end{lstlisting}

如果第一条 load 很慢，而后两条加法与它无关，理想情况下 CPU 不应该傻等。乱序执行的目标就是：在不改变架构可见结果的前提下，先执行已经 ready 的指令。

现代乱序 CPU 大致会：

\begin{itemize}
  \item 按程序顺序取指和解码。
  \item 通过 register renaming 消除一部分假依赖。
  \item 把 uops 放入调度结构。
  \item 哪些操作数 ready，哪些执行单元可用，就先执行哪些。
  \item 最后按程序顺序 retire，保证异常和结果可见性正确。
\end{itemize}

这解释了一个重要现象：你在汇编里看到的是程序顺序，但 CPU 内部的实际执行顺序可能不同。性能分析不能只看文本顺序，还要看依赖和资源。

\subsection{寄存器重命名：为什么架构寄存器不等于物理寄存器}

ISA 只暴露 16 个通用寄存器，但现代 CPU 内部通常有更多物理寄存器。寄存器重命名会把架构寄存器名映射到物理寄存器。

看这个例子：

\begin{lstlisting}
mov eax, edi
add eax, esi
mov eax, r8d
add eax, r9d
\end{lstlisting}

从文本看，四条指令都写 \register{eax}。但第二组计算并不需要第一组计算的结果。没有重命名时，对同一个架构寄存器的写会制造不必要的顺序限制。重命名可以让两个不同生命期的 \register{eax} 映射到不同物理寄存器，从而消除假依赖。

这对 C++ 优化有一个启发：不要用源代码变量名想当然地判断寄存器压力和依赖。编译器和 CPU 都会重命名、合并、消除和重排中间值。真正可靠的方式是看编译后的汇编，再结合微架构模型分析。

\subsection{Cache：为什么内存访问不是一个速度}

程序员经常说``访问内存''，但 CPU 看到的是层级结构：

\begin{center}
\begin{tabular}{lll}
\toprule
层级 & 大致特点 & 学习重点 \\
\midrule
寄存器 & 最靠近执行单元，数量少 & 分配、依赖、寄存器压力 \\
L1 data cache & 很小但很快 & cache line、局部性、load/store 吞吐 \\
L2 cache & 更大但更慢 & 工作集大小、miss 代价 \\
L3 cache & 多核心共享或分片 & 跨核心影响、容量 \\
主内存 & 容量大但延迟高 & bandwidth、NUMA、预取 \\
\bottomrule
\end{tabular}
\end{center}

CPU 不是按单个 \code{int} 从内存取数据，而是按 cache line 在 cache 和更低层级之间移动数据。常见 cache line 大小是 64 字节。顺序扫描数组通常比随机访问快，一个重要原因就是顺序访问能利用 cache line 和硬件预取。

本章只建立直觉：内存访问速度取决于数据在哪一级，以及访问模式能否被 cache 和 prefetcher 利用。后面内存层级部分会用实验系统展开。

\subsection{分支预测：CPU 为什么要猜}

流水线越深、前端越宽，分支的代价越高。如果 CPU 不猜，它遇到条件跳转就必须等条件计算完成，然后才知道下一条从哪里取。这会让前端经常空转。

branch predictor 试图预测：

\begin{itemize}
  \item 条件分支会不会跳。
  \item 跳转目标在哪里。
  \item \instruction{ret} 会返回到哪里。
\end{itemize}

如果预测正确，流水线保持忙碌。如果预测错误，CPU 必须丢弃错误路径上的推测工作，回到正确路径。这会造成明显性能损失。

这解释了为什么数据相关、难预测的分支会慢，也解释了为什么有时编译器或手写优化会使用条件移动、mask、SIMD blend 等方式减少分支。不是所有分支都坏，坏的是高频、难预测、代价明显的分支。

\subsection{从 C++ 到 CPU 行为的阅读方法}

以后看到一段 C++ hot loop，不要只停留在源代码层。建议按下面步骤阅读：

\begin{enumerate}
  \item 写出这段代码的输入、输出和允许的 C++ 语义。
  \item 编译成汇编，分别看 \code{-O0} 和 \code{-O3}，但以优化版本为性能分析对象。
  \item 标出每条关键指令读写的寄存器和内存。
  \item 找出循环携带依赖，例如 reduction 中的累加变量。
  \item 找出内存访问模式，例如连续、跨步、随机、gather。
  \item 判断控制流是否有难预测分支。
  \item 用工具验证，不靠纯想象。
\end{enumerate}

本书会不断训练这套方法。它看起来繁琐，但熟练以后会变成直觉。真正高手读性能代码时，脑子里会自动浮现数据路径、指令路径和瓶颈假设。

\subsection{贯穿例子：求和循环}

考虑：

\begin{lstlisting}[language=C++]
extern "C" int sum_i32(const int* a, long n)
{
    int s = 0;
    for (long i = 0; i < n; ++i) {
        s += a[i];
    }
    return s;
}
\end{lstlisting}

一个标量汇编形态可能包含：

\begin{lstlisting}
xor eax, eax
xor ecx, ecx
loop:
add eax, dword ptr [rdi + rcx*4]
inc rcx
cmp rcx, rsi
jl loop
ret
\end{lstlisting}

逐层理解：

\begin{itemize}
  \item \register{rdi} 保存数组起始地址。
  \item \register{rsi} 保存长度 \code{n}。
  \item \register{eax} 保存累加和。
  \item \register{rcx} 保存循环下标。
  \item \code{[rdi + rcx*4]} 是数组元素地址。
  \item \instruction{add eax, dword ptr [...]} 同时做 load 和加法。
  \item \instruction{cmp} 和 \instruction{jl} 控制循环。
\end{itemize}

性能上，这段循环至少有几个问题可以讨论：

\begin{itemize}
  \item 累加变量 \register{eax} 形成循环携带依赖。
  \item 每次迭代有一次 load。
  \item 每次迭代有循环分支。
  \item 如果编译器能证明安全并且目标支持 SIMD，优化版本可能向量化。
\end{itemize}

这就是从源代码到性能模型的第一步。你不需要在本章就优化它，但必须能把它拆成寄存器、内存、控制流和依赖。

\subsection{实验 2.1：观察机器码字节}

\begin{labbox}
\textbf{目标：}确认 CPU 执行的是机器码字节，汇编只是人类可读表示。

\textbf{步骤：}

\begin{lstlisting}[language=bash]
mkdir -p scratch/ch02
cat > scratch/ch02/add.cpp <<'EOF'
extern "C" int add_i32(int x, int y)
{
    return x + y;
}
EOF

g++ -O0 -c scratch/ch02/add.cpp -o scratch/ch02/add_O0.o
g++ -O3 -c scratch/ch02/add.cpp -o scratch/ch02/add_O3.o

objdump -drwC -Mintel scratch/ch02/add_O0.o
objdump -drwC -Mintel scratch/ch02/add_O3.o
\end{lstlisting}

\textbf{交付物：}

\begin{itemize}
  \item 截取 \code{add_i32} 的反汇编。
  \item 标出每条指令前面的机器码字节。
  \item 解释 \code{-O0} 和 \code{-O3} 的差异。
  \item 说明返回值为什么在 \register{eax}。
\end{itemize}
\end{labbox}

\subsection{实验 2.2：用 GDB 单步观察寄存器}

\begin{labbox}
\textbf{目标：}用调试器观察 \register{rip}、\register{rsp}、\register{eax}、\register{edi}、\register{esi} 如何变化。

\textbf{准备代码：}

\begin{lstlisting}[language=C++]
#include <cstdio>

extern "C" __attribute__((noinline))
int add_i32(int x, int y)
{
    return x + y;
}

int main()
{
    int z = add_i32(10, 32);
    std::printf("%d\n", z);
}
\end{lstlisting}

\textbf{编译：}

\begin{lstlisting}[language=bash]
g++ -O0 -g scratch/ch02/debug_add.cpp -o scratch/ch02/debug_add
gdb scratch/ch02/debug_add
\end{lstlisting}

\textbf{GDB 命令：}

\begin{lstlisting}
set disassembly-flavor intel
break add_i32
run
disassemble /r add_i32
info registers rip rsp rax eax rdi edi rsi esi
si
info registers rip rsp rax eax rdi edi rsi esi
si
info registers rip rsp rax eax rdi edi rsi esi
\end{lstlisting}

\textbf{交付物：}

\begin{itemize}
  \item 每次 \code{si} 后记录关键寄存器。
  \item 解释 \register{rip} 为什么变化。
  \item 解释 \register{rsp} 在函数入口和返回附近怎样变化。
  \item 解释调试版本为什么比优化版本多很多栈帧相关指令。
\end{itemize}
\end{labbox}

\subsection{实验 2.3：flags 和条件跳转}

\begin{labbox}
\textbf{目标：}理解 \instruction{cmp}、\instruction{test}、条件跳转如何配合。

\textbf{任务：}写三个函数：

\begin{lstlisting}[language=C++]
extern "C" int is_zero(int x)
{
    return x == 0;
}

extern "C" int max_i32(int a, int b)
{
    return a > b ? a : b;
}

extern "C" int abs_i32(int x)
{
    return x < 0 ? -x : x;
}
\end{lstlisting}

分别使用 \code{-O0}、\code{-O2}、\code{-O3} 编译，观察是否出现：

\begin{itemize}
  \item \instruction{cmp}
  \item \instruction{test}
  \item \instruction{jcc}
  \item \instruction{cmov}
  \item \instruction{setcc}
\end{itemize}

\textbf{交付物：}

\begin{itemize}
  \item 为每个函数画出控制流或数据流。
  \item 解释哪条指令设置 flags，哪条指令读取 flags。
  \item 判断优化版本是否减少了分支。
\end{itemize}
\end{labbox}

\subsection{实验 2.4：call、ret 和栈}

\begin{labbox}
\textbf{目标：}理解函数调用如何使用栈保存返回地址。

\textbf{任务：}写三个 noinline 函数：

\begin{lstlisting}[language=C++]
extern "C" __attribute__((noinline)) int f3(int x)
{
    return x + 3;
}

extern "C" __attribute__((noinline)) int f2(int x)
{
    return f3(x + 2);
}

extern "C" __attribute__((noinline)) int f1(int x)
{
    return f2(x + 1);
}
\end{lstlisting}

用 GDB 在 \code{f3} 里停住，执行：

\begin{lstlisting}
bt
info registers rip rsp rbp
x/8gx $rsp
disassemble /r f1
disassemble /r f2
disassemble /r f3
\end{lstlisting}

\textbf{交付物：}

\begin{itemize}
  \item 画出调用链。
  \item 标出每一层返回地址。
  \item 解释 \instruction{call} 和 \instruction{ret} 怎样配合。
  \item 比较 \code{-O0} 和 \code{-O2} 下调用链是否变化，解释 inline 的影响。
\end{itemize}
\end{labbox}

\subsection{实验 2.5：顺序访问和随机访问的第一印象}

\begin{labbox}
\textbf{目标：}用简单实验感受 cache 和访问模式对性能的影响。本实验不要求建立完整 cache 模型，只要求形成问题意识。

\textbf{任务：}分配一个较大的 \code{std::vector<int>}，分别做：

\begin{itemize}
  \item 顺序求和。
  \item 固定 stride 求和，例如 stride = 16。
  \item 按一个打乱后的 index 数组随机求和。
\end{itemize}

要求：

\begin{itemize}
  \item 用相同数据规模。
  \item 防止编译器把循环优化掉。
  \item 每个 case 至少运行多次，报告中位数。
  \item 记录 CPU 型号和编译命令。
\end{itemize}

\textbf{交付物：}

\begin{itemize}
  \item 三种访问模式的时间。
  \item 你对差异的解释。
  \item 你认为还需要哪些证据才能证明解释是正确的。
\end{itemize}
\end{labbox}

\subsection{常见误区}

\begin{warningbox}
\textbf{误区一：CPU 执行的是 C++。}

不准确。CPU 执行机器码字节。C++ 是高层语言，编译器把它翻译到目标 ISA。调试信息可以帮助你把机器状态映射回源代码，但那不是 CPU 的原生输入。
\end{warningbox}

\begin{warningbox}
\textbf{误区二：汇编和机器码是一回事。}

不准确。汇编是文本表示，机器码是字节。它们可以互相转换，但反汇编不一定恢复原始汇编源文件中的标签、注释和宏结构。
\end{warningbox}

\begin{warningbox}
\textbf{误区三：一条汇编一定对应一个硬件动作。}

不准确。复杂 x86 指令可能解码成多个 uops；简单指令也可能在流水线中经历多个阶段。性能上要区分 ISA 指令、uops、执行端口、延迟和吞吐。
\end{warningbox}

\begin{warningbox}
\textbf{误区四：程序顺序就是内部执行顺序。}

不准确。现代 CPU 可以乱序执行 ready 的操作，但 retire 时必须保持架构可见结果符合程序顺序。性能分析要看依赖图，而不只是文本顺序。
\end{warningbox}

\begin{warningbox}
\textbf{误区五：内存访问速度固定。}

不准确。同样是一条 load，数据在 L1、L2、L3、主内存时成本完全不同。访问模式、工作集大小、TLB、预取和 NUMA 都会影响性能。
\end{warningbox}

\subsection{本章作业}

\begin{exercisebox}
\textbf{作业 1：手工执行指令。}

给定函数入口状态：

\begin{itemize}
  \item \register{edi} = 7。
  \item \register{esi} = 5。
  \item \register{edx} = 3。
  \item \register{rsp} = \code{0x7fffffffe000}。
  \item memory[\code{0x7fffffffe000}] 保存返回地址 \code{0x401234}。
\end{itemize}

手工执行：

\begin{lstlisting}
mov eax, edi
add eax, esi
sub eax, edx
ret
\end{lstlisting}

写出每条指令后的 \register{eax}、\register{rax}、\register{rsp}、\register{rip} 的变化，并说明哪条指令会更新 flags。
\end{exercisebox}

\begin{exercisebox}
\textbf{作业 2：解释一个数组访问。}

对下面函数：

\begin{lstlisting}[language=C++]
extern "C" int load2(const int* a, long i)
{
    return a[i + 2];
}
\end{lstlisting}

生成优化汇编。解释地址表达式中 base、index、scale、displacement 分别是什么。说明为什么 \code{int} 下标会乘以 4。
\end{exercisebox}

\begin{exercisebox}
\textbf{作业 3：比较 \instruction{add} 和 \instruction{lea}。}

写两个或更多简单加法函数，尝试让编译器生成 \instruction{add} 或 \instruction{lea}。报告中必须包含：

\begin{itemize}
  \item 源码。
  \item 编译命令。
  \item 反汇编。
  \item 你对编译器选择的解释。
  \item 这两类指令对 flags 的影响差异。
\end{itemize}
\end{exercisebox}

\begin{exercisebox}
\textbf{作业 4：调用链报告。}

基于实验 2.4，写一份不少于 800 字的报告，回答：

\begin{itemize}
  \item 函数调用前参数在哪里？
  \item 返回值在哪里？
  \item 返回地址在哪里？
  \item \register{rsp} 如何变化？
  \item 优化级别改变后，为什么调用链可能消失？
\end{itemize}
\end{exercisebox}

\begin{exercisebox}
\textbf{作业 5：第一次性能解释。}

完成实验 2.5，写一份性能报告。报告不能只写``随机访问慢''，必须包含：

\begin{itemize}
  \item 数据规模。
  \item 编译器版本和命令。
  \item CPU 型号。
  \item 每种访问模式的时间统计。
  \item 你基于 cache line 和局部性的解释。
  \item 你认为本实验还有哪些不严谨之处。
\end{itemize}
\end{exercisebox}

\subsection{验收标准}

学完本章，你应该能够做到：

\begin{itemize}
  \item 解释 CPU 不直接执行 C++，而是执行机器码字节。
  \item 区分汇编文本、机器码、ISA 语义和微架构实现。
  \item 说清楚 \register{rip}、通用寄存器、\register{rsp}、RFLAGS 的基本作用。
  \item 手工执行简单整数指令，追踪寄存器和 \register{rip} 变化。
  \item 解释 \instruction{call} 和 \instruction{ret} 如何用栈保存和恢复控制流。
  \item 解释一条 load 指令为什么可能涉及地址生成、TLB 和 cache。
  \item 用 \code{objdump} 观察机器码字节和反汇编。
  \item 用 GDB 单步观察寄存器变化。
  \item 初步解释 pipeline、cache、branch prediction、out-of-order execution 为什么存在。
\end{itemize}

如果你能完成这些目标，后面正式进入 x86-64 汇编、ABI、性能测量和微架构时，就不会把术语当成孤立概念，而能把它们放回真实执行链路中理解。
